\subsection{Magnetic fields}

All DC magnetic fields are produced by Helmholtz coil pairs wound with copper wire on
phenolic bobbins. Three independent coil pairs are used: the main horizontal field coil,
the horizontal sweep field coil, and the vertical compensating field coil.
 
\subsubsection*{Helmholtz Coil Field}
 
A Helmholtz pair consists of two identical circular coils of radius $R$, each with $N$
turns, separated by a distance equal to $R$. The on-axis field at the midpoint between
the two coils is obtained by summing the Biot-Savart contributions from each coil.
For a single circular loop of radius $R$ carrying current $I$, the axial field at
distance $x$ from the centre is:
\begin{equation}
    B_x = \frac{\mu_0 I R^2}{2(R^2 + x^2)^{3/2}}
\end{equation}
For the Helmholtz configuration, the two coils are placed at $x = \pm R/2$. The total
field at the midpoint ($x = 0$) is:
\begin{equation}
    B = 2 \times \frac{\mu_0 N I R^2}{2\left(R^2 + (R/2)^2\right)^{3/2}}
      = \frac{\mu_0 N I R^2}{\left(\frac{5}{4}R^2\right)^{3/2}}
      = \frac{\mu_0 N I}{R}\left(\frac{4}{5}\right)^{3/2}
\end{equation}
Evaluating the numerical prefactor:
\begin{equation}
    B = \frac{8}{\sqrt{125}}\,\frac{\mu_0 N I}{R}
      \approx 0.7155\,\frac{\mu_0 N I}{R}
\end{equation}
In Gaussian units, this becomes:
\begin{equation}
    B\ (\text{gauss}) = 8.991\times10^{-3}\,\frac{NI}{R}
\end{equation}
where $I$ is in amperes and $R$ is in metres. The uniformity of the Helmholtz field is
guaranteed by the vanishing of the first and second derivatives of $B$ at the midpoint:
\begin{equation}
    \frac{\partial B}{\partial x}\bigg|_{x=0} = 0,
    \qquad
    \frac{\partial^2 B}{\partial x^2}\bigg|_{x=0} = 0
\end{equation}
so the field is uniform to second order in $x/R$.
 
\subsubsection*{Coil Specifications}
 
The properties of the three coil pairs are summarised in Table~\ref{tab:coils}.
 
\begin{table}[h]
    \centering
    \begin{tabular}{lcccc}
        \hline
        Coil & Mean radius & Turns/side & Field/Amp & Max.\ field \\
             & (cm)        &            & ($10^{-4}$~T/A) & ($10^{-4}$~T) \\
        \hline
        Vertical   & 11.735 & 20  & 1.5 & 1.5 \\
        Horizontal & 15.790 & 154 & 8.8 & 22.0 \\
        Sweep      & 16.390 & 11  & 0.60 & 0.60 \\
        \hline
    \end{tabular}
    \caption{Magnetic field coil specifications.}
\end{table}
 
\subsubsection*{Current-Regulated Field Control}
 
Each coil pair is driven by a voltage-to-current converter circuit. The reference
voltage $V_{\text{ref}}$ sets the voltage across a sense resistor $R_s$, thereby
fixing the current through the coil:
\begin{equation}
    I_{\text{coil}} = \frac{V_{\text{ref}}}{R_s}
\end{equation}
The magnetic field is then:
\begin{equation}
    B = \frac{8}{\sqrt{125}}\,\frac{\mu_0 N}{R} \cdot \frac{V_{\text{ref}}}{R_s}
\end{equation}
The sense resistors for the three coil pairs are:
\begin{equation}
    R_s^{\text{vertical}} = 1\ \Omega, \qquad
    R_s^{\text{horizontal}} = 0.5\ \Omega, \qquad
    R_s^{\text{sweep}} = 1\ \Omega
\end{equation}
The voltage across $R_s$ is monitored via front-panel tip jacks and used to determine
the coil current accurately.
 
\subsubsection*{Horizontal Main Field}
 
The main horizontal field coil ($R = 15.79$~cm, $N = 154$ turns/side) produces the
primary static field along the optical axis. With the internal power supply, the maximum
current is $\approx 1.0$~A, giving:
\begin{equation}
    B_{\text{max}} = 8.8\times10^{-4}\ \text{T/A} \times 1.0\ \text{A}
                   = 8.8\times10^{-4}\ \text{T} = 8.8\ \text{gauss}
\end{equation}
When an external power supply is used (max.\ 40~V, 3.0~A), the coil resistance at room
temperature is $R_{\text{coil}} \approx 10\ \Omega$, rising to $\approx 12\ \Omega$ at
high current due to Joule heating. The temperature coefficient of resistance for copper
gives:
\begin{equation}
    R(T) = R_0\left[1 + \alpha_{\text{Cu}}(T - T_0)\right],
    \qquad \alpha_{\text{Cu}} = 3.9\times10^{-3}\ \text{K}^{-1}
\end{equation}
The fractional change in resistance (and hence in the calibration factor) over a
temperature rise of $\Delta T \approx 50$~K is:
\begin{equation}
    \frac{\Delta R}{R_0} = \alpha_{\text{Cu}}\,\Delta T
    = 3.9\times10^{-3} \times 50 \approx 19.5\%
\end{equation}
This must be accounted for when high-current calibrations are performed.
 
\subsubsection*{Horizontal Sweep Field}
 
The sweep field coil ($R = 16.39$~cm, $N = 11$ turns/side) is wound on top of the
main horizontal coil. Its reference voltage is the sum of three contributions:
\begin{equation}
    V_{\text{ref}}^{\text{sweep}} = V_{\text{start}} + V_{\text{ramp}}(t) + V_{\text{mod}}
\end{equation}
The ramp voltage $V_{\text{ramp}}(t)$ increases linearly from 0 to $V_{\text{range}}$
over a sweep time $t_s \in [1,\,1000]$~s:
\begin{equation}
    V_{\text{ramp}}(t) = \frac{V_{\text{range}}}{t_s}\,t, \qquad 0 \leq t \leq t_s
\end{equation}
The corresponding field sweep rate is:
\begin{equation}
    \dot{B} = \frac{8}{\sqrt{125}}\,\frac{\mu_0 N}{R\,R_s}\cdot
              \frac{V_{\text{range}}}{t_s}
\end{equation}
The modulation voltage $V_{\text{mod}}$ allows a small sinusoidal or square-wave
modulation of the field for lock-in detection:
\begin{equation}
    V_{\text{mod}} = A_{\text{mod}}\,f(t), \qquad
    \frac{V_{\text{mod}}}{B_{\text{mod}}} \approx 1\ \text{V per}\ 10\ \text{mG}
\end{equation}
 
\subsubsection*{Vertical Compensating Field}
 
The vertical coil ($R = 11.735$~cm, $N = 20$ turns/side) compensates for the vertical
component of the Earth's magnetic field $B_{\text{Earth}}^{\perp}$. In Buffalo, NY
(geomagnetic latitude $\approx 53°$), the vertical component is approximately:
\begin{equation}
    B_{\text{Earth}}^{\perp} = B_{\text{Earth}}\sin\lambda_m
    \approx 5\times10^{-5}\ \text{T} \times \sin 53°
    \approx 4.0\times10^{-5}\ \text{T} = 0.40\ \text{gauss}
\end{equation}
The required compensating current is:
\begin{equation}
    I_{\text{comp}} = \frac{B_{\text{Earth}}^{\perp}}{\left(8/\sqrt{125}\right)\mu_0 N/R}
    \approx 0.33\ \text{A}
\end{equation}
The vertical coil is wired so that the field points downward when the red plug is in the
red jack, which is the correct polarity to cancel the Earth's vertical field in the
Northern Hemisphere.
 
\subsubsection*{Recorder Output}
 
The recorder output provides a signal proportional to the sweep coil current, derived
from the voltage across the $1\ \Omega$ sense resistor and amplified with a low-pass
filter of time constant $\tau = 2$~ms. The conversion factor is:
\begin{equation}
    \frac{V_{\text{rec}}}{B_{\text{sweep}}} = 50\ \text{mV per}\ 1\ \text{mG}\ (10\ \mu\text{T})
\end{equation}
with an output range of $-13.5$~V to $+13.5$~V.